\chapter{Preliminaries}\label{ch:prelim}
\section{Notation and conventions}\label{sec:prelim}

Throughout, $P=(p_0,\dots,p_{n-1})$ is a sequence of $n\ge3$ points in
\emph{strictly} convex position (no three collinear), listed counterclockwise
(ccw); indices are taken modulo~$n$. We write $|pq|$ for the Euclidean
distance. For $i\ne j$, $P(i,j)$ denotes the points strictly between $p_i$ and
$p_j$ when going ccw from $p_i$ to $p_j$. For $r>0$, a set is \emph{$r$-far}
if all its pairwise distances are at least $r$, and $\far_r(i,j)$ is the
predicate $|p_ip_j|\ge r$. For three points, $\ang abc\in(0,\pi)$ is the angle
at $b$ in the triangle $abc$, and $\circlethree abc$ is the circle through
them. $\per(P)$ is the perimeter of $\conv(P)$.

A \emph{triangulation} of a set $S$ in convex position, $|S|\ge3$, is a
maximal set of pairwise non-crossing segments (\emph{edges}) between points of
$S$; its edges are the sides of $\conv(S)$ and the \emph{diagonals}. A
diagonal $ab$ with incident triangles $abc$ and $bad$ is \emph{locally
Delaunay} if $d$ is not strictly inside $\circlethree abc$. Since $acbd$ is a
convex quadrilateral, this is equivalent to
\begin{equation}\label{eq:ld}
\ang acb+\ang adb\le\pi .
\end{equation}
A triangulation is \emph{Delaunay} if no point of $S$ lies strictly inside the
circumcircle of any of its triangles. Delaunay triangulations exist for every
finite point set not contained in a line (triangulate the cells of the
Delaunay subdivision arbitrarily).

We repeatedly use two elementary facts.
\begin{description}[leftmargin=1.5em,itemsep=1pt]
\item[(Q)] In a convex quadrilateral $abcd$ the diagonals are together longer
than either pair of opposite sides: $|ac|+|bd|>|ab|+|cd|$ and
$|ac|+|bd|>|bc|+|da|$. (Apply the triangle inequality at the crossing point of
the diagonals.)
\item[(A)] Seen from a vertex $v$ of a convex polygon, the other vertices
appear in their cyclic order within an angle smaller than $\pi$. Hence, if
$v,u_1,u_2,u_3$ occur in this ccw order, $\ang u_1vu_3=\ang u_1vu_2+\ang u_2vu_3$.
\end{description}

\paragraph{Model of computation and exact predicates.}
Our algorithms compare distances and evaluate the in-circle predicate; no other
geometric primitive is used. The analysis holds on a real RAM, where the term
$n^3/w$ below is to be read as $n^3$, and on a word RAM with word size
$w=\Omega(\log n)$ for integer input. For integer coordinates of absolute value
at most $2^{29}$, squared distances fit into $62$ bits and the in-circle
determinant is smaller than $2^{124}$ in absolute value, so that signed
$128$-bit integers suffice; this is what our implementation does. All angle
comparisons below reduce to in-circle tests by the inscribed angle theorem:
\begin{description}[leftmargin=1.5em,itemsep=1pt]
\item[(I1)] if $x$ and $y$ lie on the same side of the line $ab$, then
$\ang axb>\ang ayb$ iff $x$ is strictly inside $\circlethree ayb$;
\item[(I2)] if $x$ and $c$ lie on opposite sides of the line $ab$, then
$\ang axb+\ang acb>\pi$ iff $x$ is strictly inside $\circlethree acb$.
\end{description}

\section{A Delaunay certificate for \texorpdfstring{$r$}{r}-far sets}\label{sec:certificate}

\begin{lemma}[Delaunay lemma {\citep{delaunay,edelsbrunner}}]\label{lem:delaunay}
If every diagonal of a triangulation $T$ of a finite point set $S$ is locally
Delaunay, then $T$ is a Delaunay triangulation of $S$.
\end{lemma}

It is well known that a closest pair of a point set is joined by an edge of
its Delaunay triangulation. We need the statement for an \emph{arbitrary}
Delaunay triangulation in the presence of cocircular points, and give the
short proof for sets in convex position.

\begin{lemma}\label{lem:closest}
Let $T$ be a Delaunay triangulation of a set $S$ in strictly convex position,
$|S|\ge3$. Then every closest pair of $S$ is joined by an edge of $T$.
\end{lemma}
\begin{proof}
Let $p,q$ be a closest pair and suppose that $pq$ is not an edge of $T$. Then
$pq$ is not a side of $\conv(S)$, and it does not pass through a third point
of $S$. Hence, starting at $p$, the segment $pq$ runs through the interior of
a triangle $puv$ of $T$ and leaves it through an interior point of the side
$uv$. So $u$ and $v$ lie strictly on opposite sides of the line $pq$, and
$p,u,q,v$ are the vertices of a convex quadrilateral in this cyclic order,
with diagonals $pq$ and $uv$. Because $T$ is Delaunay, $q$ is not strictly
inside $\circlethree puv$, that is, $\ang upv+\ang uqv\le\pi$ by~(I2). The
angles of the quadrilateral sum to $2\pi$, so $\ang puq+\ang pvq\ge\pi$, and
one of them, say $\ang puq$, is at least $\pi/2$. Then $pq$ is the strictly
longest side of the triangle $puq$, so $|pu|<|pq|$, contradicting the choice
of $p,q$.
\end{proof}

\begin{theorem}[Delaunay certificate]\label{thm:cert-delaunay}
Let $S$ be a set in strictly convex position, $|S|\ge3$, and let $r>0$. Then
$S$ is $r$-far if and only if $S$ has a triangulation $T$ such that
\begin{enumerate}[label=(\alph*),itemsep=0pt,topsep=2pt]
\item every edge of $T$ has length at least $r$, and
\item every diagonal of $T$ is locally Delaunay.
\end{enumerate}
\end{theorem}
\begin{proof}
($\Leftarrow$) By (b) and \cref{lem:delaunay}, $T$ is a Delaunay
triangulation. By \cref{lem:closest} the minimum pairwise distance of $S$ is
the length of an edge of $T$, which is at least $r$ by~(a).
($\Rightarrow$) Take any Delaunay triangulation $T$ of $S$. It satisfies (b)
by definition, and (a) because all its edges join points of the $r$-far
set~$S$.
\end{proof}

Conditions (a) and (b) are \emph{local}: (a) concerns single edges and (b)
pairs of adjacent triangles. The dynamic program of
Tkachenko and Wang~\citep{tw26} is built on the same idea of gluing triangles
under a local compatibility condition (for disks it is phrased with
Voronoi vertices of triples); \cref{thm:cert-delaunay} isolates the
statement that makes such a program correct and reduces its proof to the
Delaunay lemma.

